\section{PAPER 027: AutonomousDiscoveryLoop — Closing the Discover-Validate-Approve-Materialize Loop for Knowledge Graph Self-Improvement}

\textbf{Author}: Bryan Alexander Buchorn  
\textbf{Affiliation}: Independent Researcher, Las Vegas, NV, USA  
\textbf{Status}: v2.51.0 (Phase 167 (STRB) COMPLETE)
\textbf{Date}: May 2, 2026

\textbf{CEREBRUM Phase 74}

---

\subsection{Abstract}

We present \textbf{AutonomousDiscoveryLoop}, the orche\-stration component that closes the end-to-end autonomous knowledge graph improvement loop in CEREBRUM. The loop runs \texttt{ResearchAgent.scan\_once()} on a confi\-gurable timer, processes each finding through \texttt{AutoApprover} (Phase 71), and applies a \textbf{circuit breaker} — a sliding-window approval-rate monitor that pauses mater\-ialization when quality degrades below a threshold. A \textbf{per-cycle cap} (\texttt{max\_materializations\_per\_cycle}) prevents runaway edge insertion during anomalous discovery bursts. A \textbf{dry\_run} mode allows safe production trials without writing any edges. \texttt{AutoApprover} state is check\-pointed to disk after every cycle, enabling warm restart. The loop exposes a structured REST API for monitoring, confi\-guration, and lifecycle management, making it suitable for embedding in production CEREBRUM deployments without operator inter\-vention.

---

\subsection{1. Archit\-ecture}

\subsubsection{1.1 Loop Lifecycle}

\begin{lstlisting}
start() → background thread
    ↓
while running:
    scan_once() → findings
    for finding in findings[:cap]:
        decide = approver.decide(finding)
        if decide == APPROVE: approve(finding)
        if decide == REJECT:  reject(finding)
    record_cycle()
    checkpoint_approver()
    check_circuit_breaker()
    sleep(next_interval)
stop() → graceful shutdown
\end{lstlisting}

\subsubsection{1.2 LoopConfig Dataclass}

\begin{lstlisting}
@dataclass
class LoopConfig:
    cycle_interval: float = 300.0
    max_materializations_per_cycle: int = 10
    min_approval_rate: float = 0.5
    circuit_breaker_window: int = 20
    dry_run: bool = False
    auto_rollback_on_trip: bool = False   # Phase 79
    adaptive_tuning: bool = False          # Phase 82
    adaptive_min_cap: int = 1
    adaptive_max_cap: int = 20
    adaptive_min_interval: float = 60.0
    adaptive_max_interval: float = 7200.0
    approver_checkpoint_path: Optional[str] = None
\end{lstlisting}

\subsubsection{1.3 CycleRecord Dataclass}

\begin{lstlisting}
@dataclass
class CycleRecord:
    cycle_number: int
    started_at: float
    findings_scanned: int
    materializations: int
    approvals: int
    rejections: int
    reviews: int
    circuit_breaker_tripped: bool
    edges_rolled_back: int = 0    # Phase 79
    effective_cap: int = 0        # Phase 82
\end{lstlisting}

---

\subsection{2. Circuit Breaker}

The circuit breaker computes approval rate over a sliding window of the last N decisions (approvals + rejections):

\begin{lstlisting}
approval_rate = approvals_in_window / (approvals + rejections)_in_window
\end{lstlisting}

If \texttt{approval\_rate \textless  min\_approval\_rate}:
\begin{itemize}
\item \texttt{circuit\_breaker\_tripped = True}
\item Materi\-alization pauses for the current cycle
\item If \texttt{auto\_rollback\_on\_trip=True} (Phase 79): \texttt{ProvenanceLedger.rollback\_cycle()} is called
\item Loop continues sleeping; next cycle attempts recovery

\end{itemize}
The circuit breaker resets when the approval rate recovers above threshold over a fresh window.

---

\subsection{3. Per-Cycle Cap}

\texttt{max\_materializations\_per\_cycle} acts as a hard upper bound on approved findings per cycle. Even if \texttt{AutoApprover} would approve 100 findings in a single scan, only the first N are mater\-ialized. This prevents:
\begin{itemize}
\item Burst mater\-ialization that overwhelms downstream systems
\item Runaway edge accum\-ulation during transient high-discovery phases

\end{itemize}
Combined with \texttt{adaptive\_tuning} (Phase 82), the cap is dynamically scaled per cycle based on \texttt{DiscoveryCalibrator} community weights.

---

\subsection{4. Dry Run Mode}

\texttt{LoopConfig(dry\_run=True)} makes the loop execute all phases (scan, validate, decide) but skips \texttt{approve()} / \texttt{reject()} calls. The loop records what \textit{would} have been mater\-ialized in \texttt{CycleRecord}, enabling safe production trials to measure expected impact before enabling writes.

---

\subsection{5. REST API}

\begin{center}
{\footnotesize
\begin{tabularx}{\columnwidth}{|>{\raggedright\arraybackslash}X|>{\raggedright\arraybackslash}X|>{\raggedright\arraybackslash}X|}
\hline
Endpoint & Method & Description \\ \hline
\texttt{/research/loop/start} & POST & Start loop (idempotent) \\ \hline
\texttt{/research/loop/stop} & POST & Graceful stop \\ \hline
\texttt{/research/loop/status} & GET & Running state, cycle history, approval rate, circuit breaker \\ \hline
\texttt{/research/loop/configure} & POST & Partial update: \texttt{cycle\_interval}, \texttt{max\_materializations\_per\_cycle}, \texttt{min\_approval\_rate}, \texttt{dry\_run}, \texttt{auto\_rollback\_on\_trip}, \texttt{adaptive\_tuning} \\ \hline
\end{tabularx}
}
\end{center}

---

\subsection{6. References}

\begin{itemize}
\item [dean\-2008\-mapreduce] Dean, J. \& Ghemawat, S. (2008). MapReduce: Simplified data processing on large clusters. \textit{CACM}, 51(1), 107–113.

\end{itemize}
---
\textbf{Copyright © 2026 Bryan Alexander Buchorn. All Rights Reserved.}

---
\subsection{Acknow\-ledgments}

The author gratefully ackno\-wledges the use of Claude (Anthropic) as a research assistant throughout this work. Claude assisted with mathe\-matical forma\-lization, code generation, manuscript preparation, and technical writing. All conceptual contr\-ibutions, archi\-tectural decisions, exper\-imental design, and intel\-lectual claims are solely the author's.

\textbf{Reviewed on}: May 2, 2026 for version v2.51.0
